% Exemplo: biologia com genealogytree (heredograma com indivíduos afetados; quadrado, círculo e
% preenchimento definidos como estilos próprios) e um quadro de Punnett feito com tabular.
% Compilar: latexmk -pdf biologia.tex             Testado com: genealogytree 2.4, TeX Live 2024 ou mais recente
%%BEGIN:biopreambulo
\documentclass[12pt, a4paper]{article}
\usepackage[T1]{fontenc}
\usepackage[brazilian]{babel}
\usepackage[all]{genealogytree}       % heredogramas e árvores genealógicas
\usepackage{array}
%%END:biopreambulo

\begin{document}
\pagestyle{empty}

%%BEGIN:bioheredograma
Heredograma de uma característica recessiva (símbolos cheios: afetados):
\begin{center}
\begin{genealogypicture}[node size=9mm, level size=9mm, level distance=9mm,
    box={colback=white, colframe=black, boxrule=0.6pt, boxsep=0pt,
         left=0pt, right=0pt, halign=center, valign=center, fontupper=\tiny},
    edges={foreground={black, line width=0.6pt}, background={white}},
    homem/.style={box={sharp corners}}, % quadrado
    mulher/.style={box={arc=4.5mm}},    % círculo (arco = metade do lado)
    afetado/.style={box={colback=black!75, colupper=white}}]
  child{
    g[homem]{I-1} p[mulher]{I-2}
    c[mulher]{II-1}
    c[homem, afetado]{II-2}
    child{
      g[mulher]{II-3} p[homem]{II-4}
      c[homem]{III-1}
      c[mulher, afetado]{III-2}
    }
  }
\end{genealogypicture}
\end{center}
%%END:bioheredograma

%%BEGIN:biopunnett
Cruzamento \textit{Aa} $\times$ \textit{Aa}:
\begin{center}
\renewcommand{\arraystretch}{1.6}
\begin{tabular}{c|>{\centering\itshape}p{1cm}|%
                >{\centering\itshape\arraybackslash}p{1cm}|}
  \multicolumn{1}{c}{} & \multicolumn{1}{c}{\textit{A}}
    & \multicolumn{1}{c}{\textit{a}} \\
  \cline{2-3}
  \textit{A} & AA & Aa \\ \cline{2-3}
  \textit{a} & Aa & aa \\ \cline{2-3}
\end{tabular}
\end{center}
Proporção genotípica 1~\textit{AA} : 2~\textit{Aa} : 1~\textit{aa}.
%%END:biopunnett

\end{document}
